\chapter*{Document Control}
\addcontentsline{toc}{chapter}{Document Control}
\markboth{Document Control}{Document Control}

\section*{Purpose of this document set}

This documentation set describes the \textbf{Expo Snips} repository located at
\RepositoryPath. It was produced by a complete reverse-engineering analysis of
the source tree: package manifests, TypeScript sources, registry data, Handlebars
templates, build scripts, GitHub Actions workflows and Markdown documentation.
No application source code was modified.

\section*{Revision history}

\begin{table}[H]
\centering
\caption{Document revision history}
\begin{tabular}{p{2cm}p{3cm}p{9.4cm}}
\toprule
\textbf{Version} & \textbf{Date} & \textbf{Change} \\
\midrule
1.0 & 17 October 2026 & Initial reverse-engineering documentation suite
(45 chapters/appendices, 10 Mermaid diagrams, PDF and DOCX outputs). \\
1.1 & 19 October 2026 & Refresh against the current tree: the three crashing
catalog templates were fixed (escaped Handlebars literals and typed component
props --- see Chapter~\ref{ch:known-issues} and the File Index), 22 committed
screen screenshots were added under \texttt{docs/screenshots/} and are now
referenced from the UI/UX chapter, the root workspace was returned to a
dependency-free state, and commit/branch figures were updated. \\
\bottomrule
\end{tabular}
\end{table}

\section*{Source of truth}

The repository is the only source of truth. Every material statement in this
document set is traceable to a file, symbol, configuration key or commit. Where
the repository does not contain evidence for a statement, the statement is either
omitted or explicitly labelled \textbf{UNKNOWN}.

\section*{Evidence classification scheme}

Each material claim in this documentation carries a status label:

\begin{itemize}
  \item \textbf{IMPLEMENTED} --- present and verifiable in the repository source.
  \item \textbf{INFERRED} --- a reasonable conclusion drawn from repository evidence,
        but not stated explicitly by the code.
  \item \textbf{RECOMMENDED} --- a suggestion made by this document; it is \emph{not}
        a description of existing behaviour.
  \item \textbf{UNKNOWN} --- no evidence could be found in the repository.
  \item \textbf{NOT APPLICABLE} --- the subject does not exist in this product
        (for example, a database, because the product has no data store).
\end{itemize}

The full scheme and a worked example are given in Appendix~C.

\section*{Scope of analysis}

\begin{table}[H]
\centering
\caption{Analysis scope and artefacts inspected}
\begin{longtable}{p{4.4cm}p{10cm}}
\toprule
\textbf{Area} & \textbf{Artefacts inspected} \\
\midrule
\endhead
Root workspace & \texttt{package.json} (no \texttt{dependencies} block --- the workspace root is dependency-free),
\texttt{pnpm-workspace.yaml}, \texttt{tsconfig.base.json}, \texttt{.gitignore} \\
Packages & \texttt{packages/core}, \texttt{packages/cli}, \texttt{packages/registry}, \texttt{packages/vscode-extension} \\
Registry data & \texttt{packages/registry/data/**} (24 recipes), \texttt{packages/registry/templates/**} (48 templates) \\
Built registry & \texttt{registry/index.json}, \texttt{registry/recipes/**}, \texttt{packages/registry/bundle.json} \\
Scripts & \texttt{scripts/build-registry.mjs}, \texttt{scripts/lib/build-recipes.mjs}, \texttt{scripts/validate.mjs}, \texttt{scripts/serve.mjs} \\
CI/CD & \texttt{.github/workflows/deploy-registry.yml}, \texttt{.github/workflows/publish.yml} \\
Documentation & \texttt{README.md} (includes a Screen Showcase gallery), \texttt{CHANGELOG.md}, \texttt{AGENTS.md}, \texttt{documents/ProjectDescription.md}, \texttt{packages/vscode-extension/README.md} \\
Screen screenshots & \texttt{docs/screenshots/} --- 22 committed PNGs (21 screens plus a gallery shot) with regeneration notes in \texttt{docs/screenshots/README.md} \\
Version control & \texttt{git log} (18 commits at analysis time), branches \texttt{master}, \texttt{update-fix}, \texttt{origin/gh-pages} and several \texttt{dependabot/*} branches \\
\bottomrule
\end{longtable}
\end{table}

\section*{Explicitly excluded}

\begin{itemize}
  \item Generated dependency trees (\texttt{node\_modules/}) were excluded, except
        where a manifest was needed to identify a dependency.
  \item Build output directories (\texttt{dist/}, \texttt{out/}, \texttt{*.vsix}) were
        treated as generated artefacts.
  \item Git-ignored generated sources (\texttt{packages/*/src/generated/}) were
        inspected only to confirm their build-time origin.
\end{itemize}

\section*{Terminology note}

The product uses ``screen recipes'' (JSON definitions) rendered through
Handlebars templates. Throughout this document, \emph{recipe} means a registry
entry such as \texttt{auth/login}, and \emph{generated screen} means the resulting
files written into a developer's project.

\clearpage
